\begin{abstract}
Compiler optimization can significantly influence program execution time, yet determining an effective ordering of optimization passes remains challenging because the impact of a pass depends on the sequence in which transformations are applied. This work investigates sequence-aware Bayesian optimization for compiler pass ordering using the LLVM compilation framework. The proposed approach uses a Gaussian Process surrogate with a sequence-aware kernel, Expected Improvement for candidate selection, and warm-start information across benchmark programs. It is compared with Random Search, a Genetic Algorithm, Tree-structured Parzen Estimator (TPE) optimization, and a variant without warm start on five PolyBench/C benchmarks, using 20 evaluations per search and five random seeds. After correctness checking and re-timing, the geometric-mean speedup relative to LLVM \texttt{-O3} is 0.942 for the proposed method, 0.952 for TPE, 0.947 for Random Search, and 0.917 for the Genetic Algorithm. Paired Wilcoxon tests show no significant difference in verified speedup between the proposed method and any other method (all $p > 0.3$), and removing the warm start does not change the result. No method matches \texttt{-O3} on average. In exploratory comparisons of search cost, the proposed method produced fewer failed builds than TPE and spent less time per search than all baselines, mainly because its candidates executed faster. Its verified performance advantage therefore remains unproven under this budget.
\end{abstract}
\begin{IEEEkeywords}
Bayesian optimization, compiler autotuning, compiler optimization, LLVM, optimization pass ordering, program performance.
\end{IEEEkeywords}
